\section{Reproducibility Statement}
\label{sec:reproducibility}

All deterministic claims are proved in the appendices.  The universal causal and incidence theorem is proved in \cref{app:theory-strengthening}; Appendix~\ref{app:interaction-proof} reproduces the attributed P1 reduction proof for self-containment and proves the new observable-readout transfer; the observation/inference theorem is proved in \cref{app:deep-theory}; and the geometric expressivity results are proved in \cref{app:proofs}.

The deterministic verification companion is \path{experiments/optimizer_calculus/verify_calibrated_calculus.py}.  It checks incidence inversion, design rank, quotient algebra, Gaussian risk and coverage, sector tests, misspecification, randomization calibration, and action certificates.  The complete factorial closure is implemented by \path{experiments/optimizer_calculus/interaction_curvature_experiment.py} with \path{experiments/optimizer_calculus/interaction_curvature_config.json}; six derivative, curvature, design, inference, and finite-response tests are in \path{experiments/optimizer_calculus/test_interaction_curvature_experiment.py}, and \path{experiments/optimizer_calculus/verify_interaction_curvature_results.py} independently audits the saved CSV/JSON artifacts.  The submission includes the configuration snapshot, all summary tables, 4500 confirmatory minibatch observations, and result JSON under \path{newpapers/03-geometric-nongeometric-optimizer-calculus/artifacts/interaction_curvature}.  Quadratic diagnostics use \path{experiments/toy/gng_optimizer_benchmark.py}.  Trace audits use the scripts under \path{experiments/lane_audit}.

All random sources, quadrature orders, budgets, bootstrap draws, and power targets for the factorial closure are fixed in the configuration.  The production outputs pass the independent artifact audit; the six new tests and the original calibrated-calculus verification pass.  Exact software-package versions are not locked, so bitwise reproduction across numerical-library versions is not claimed.

The long runs are \path{paper03_rich_family_b1024_l28_20260710_2223} and \path{paper03_rich_cifar10_b1024_l28_20260710_2223}.  They use batch size $1024$, lane count $28$, and $4096$ recorded steps per job.  MNIST/Fashion-MNIST uses seeds $7,8,9$; CIFAR-10 uses $7,8,9,11$.  Follow-up runs are \path{paper03_followup_sector_cifar_b1024_l28_20260711_0242} and \path{paper03_followup_tinyresnet_cifar_b1024_l28_20260711_0242}, using seeds $7,8,9$ and $128$ or $96$ recorded steps.  Analysis commands, recorded parameters, job counts, model widths, splits, and monitor statistics are in \cref{app:diagnostic-details}.  The surviving run artifacts do not bind exact Python, NumPy, PyTorch, and CUDA versions to these Paper03 jobs.

The neural scripts fix Python, NumPy, PyTorch, and data-loader seeds, but their archived configuration does not assert deterministic CUDA algorithms across hardware.  The follow-up raw run directories and per-seed summaries are present in the research workspace; the long-run raw trace archives are absent from the submission, which retains their aggregate CSV/JSON summaries.  The run manifest records the launch argument vector but contains \texttt{git\_error} in place of a commit hash.  There is no environment lockfile or complete hyperparameter-search log.  Consequently the controlled factorial closure is independently auditable from supplied artifacts, while exact bitwise regeneration of every legacy neural trace is not guaranteed.

\section{Ethics Statement}
\label{sec:ethics}

This work studies optimizer theory and diagnostic experiments and uses no human-subject data.  The main risk is overstating causal attribution or optimizer performance from passive or incomplete interventions.  We separate controlled factorial effects and support tests from passive response-class exclusion and from still-unmeasured joint effects in the neural audits.
